\section{Conclusion}
\label{sec:conclusion}

We presented \method{}, a framework built around a language-grounded risk
concept interface for accident anticipation. The central contribution of the
paper is to show that such an interface can be constructed explicitly and
studied as a scientific object. In our setting, that means treating ontology
construction as part of the method rather than as hidden preprocessing, then
training an interpretable anticipation model on top of the resulting ontology.

This shift in perspective changes what the paper contributes. First, we provide
a reproducible pipeline for mining, refining, merging, balancing, and lightly
reviewing risk concepts from driving frames. Second, we show that the resulting
ontology is not merely descriptive: when deployed as the concept bottleneck of
the model, ontology choice changes the AP--mTTA operating point under matched
launchers. Third, we show that strong anticipation can coexist with a compact
and auditable semantic layer, including a competitive interpretable tier on DAD
and cleaner timing behavior on A3D.

Empirically, the paper is strongest on the semantic layer: the ontology
pipeline, the concept bottleneck, the controlled ontology comparisons, and the
case-study and audit artifacts. The timing module remains important as the
downstream consumer of semantic state, but the paper should be read first as a
semantic-interface contribution and only second as a policy-timing paper. The
actor-policy results are therefore useful mainly insofar as they clarify how
much of that behavior transfers beyond the classifier trajectory.

\paragraph{Limitations and future work.}
The concept labels are still based on CLIP-style pseudo-annotations rather than
human-verified frame-level concept supervision. The current semantic interface
is therefore auditable and trainable, but not yet a finished concept benchmark.
The evaluation also relies on offline visual features rather than end-to-end
video backbones. More importantly, while the timing module consumes the concept
state by construction, a complete policy-level intervention validation study
remains future work. We have strengthened the language-facing evidence with a
larger frame audit and full light-touch review of the released ontology, so the
main remaining empirical weakness is no longer missing ontology coverage but
DAD-side instability, limited cross-dataset ontology transfer evidence, and the
gap between classifier-derived anticipation scores and mature policy-level
timing behavior.

\paragraph{Evidence scope.}
To keep methodological integrity, the paper explicitly treats DAD as a
stress setting, not a universal performance frontier, and treats actor-policy
timing as a partial, support-only account. This choice keeps the central claim
on the part of the evidence that is strongest: concept-interface construction,
matched protocol contrast, and clean A3D flagship stability.

\paragraph{Broader impact.}
For safety-critical prediction, the practical value of \method{} is not merely
that it forecasts accidents, but that it exposes a more inspectable semantic
state than a purely opaque predictor. This could help engineers reason about
failure modes, ontology design, and warning behavior. At the same time, the
current paper does not establish that exposed concepts or CRS scores are
already reliable enough for direct human override in a deployment pipeline.
Any real-world use would still require much stronger calibration, robustness,
and human-factors validation.

\paragraph{Final takeaway.}
\method{} is strongest as a semantic-interface methodology paper:
ontology construction is explicit, auditable, and empirically linked to model
behavior.
The main remaining research direction is a tighter, uniformly decisive stress
envelope for short-horizon timing on DAD; ontology coverage is no longer the
primary gap.
